%% Example GayaUKM thesis in English
\documentclass[english,nohyphen]{GayaUKM}

\usepackage{graphicx}
\usepackage{longtable} %% If your list of symbols, abbreviations is longer than a page and you want to use a longtable for it

%% Your metadata. Note that an English thesis needs a Malay title page as well,
%% so you'll need to specify the Malay title, faculty and degreetype.
\title{A Non Invasive Approach To Depression Classification in Older Adults Leveraging Facial Expression Recognition}
\titlems{Pendekatan Bukan Invasif Terhadap Klasifikasi Kemurungan pada Orang Dewasa yang Lebih Tua Memanfaatkan Pengecaman Ekspresi Wajah}
\author{Muhammad Daffa Zahrandika Wibisono}  %% Assuming your name is spelt the same way in English and Malay
\authorid{P147483}
\faculty{Institute of Visual Informatics}
\facultyms{Institute Visual Informatik}
\submissiondate{07 April 2026}
\submissionyear{2026}
\degreetype{Master by Research}
\degreetypems{Sarjana Sains}
\campus{Bangi} %% Assuming Malaysian cities are spelt the same way in English and Malay

%% If you find the boxes around hyperlinks distracting
\hypersetup{colorlinks,allcolors=black}

\begin{document}

%% Generate the cover page
\makecoverpage

%% Generate the English and Malay title pages.
%% Re-specify your title with different manual line breaks for the English
%% title page, if necessary
\title{A Non Invasive Approach To Depression Classification in Older Adults Leveraging Facial Expression Recognition}
\maketitlepage

%%%% If you're happy to just use the same line-breaking scheme for the English
%%%% title on both the cover and the titlepage, then you can just call
%%%% \maketitle which combines both \makecoverpage and \maketitlepage.

\frontmatter
\declaration

% Acknolwedgements from ack.tex
\input{ack}

% English abstract from abstract-en.tex
\input{abstract-en}

% Malay Abstract from abstrak-ms.tex
% The Malay translation of your title needs to be given here
\begin{msAbstract}[Pendekatan Bukan Invasif Terhadap Klasifikasi Kemurungan pada Orang Dewasa yang Lebih Tua Memanfaatkan Pengecaman Ekspresi Wajah]
Inilah abstrak dalam Bahasa Melayu. Data korpus merupakan data bahasa
Melayu yang datangnya dalam dua bentuk sumber, iaitu bentuk tulisan
dan bentuk lisan. Bentuk tulisan seperti buku, majalah, surat khabar,
makalah, monograf, dokumen, kertas kerja, efemeral, puisi, drama,
kad bahan, surat, risalah dan sebagainya. Sementara bentuk lisan yang
ditranskripsikan seperti ucapan, wawancara, temu bual, perbualan dan
sebagainya dalam pelbagai bentuk rakaman.
\end{msAbstract}


\tableofcontents
\listoffigures
\listoftables

% List of Symbols may be prepared as in symbols.tex
\input{symbols}


\mainmatter
% Each chapter from a separate file
\chapter{Introduction}
\label{chap:intro}

\section{Research Background}
A significant demographic transition has occurred, marked by unprecedented demographic shifts, where increasing life expectancy and declining fertility rates are triggering global population aging. The Southeast Asian (ASEAN) region is one of the main epicenters of this transition. According to the \cite{ASEANstats2023} report, the proportion of the region's elderly population experienced a significant escalation from 5.3\% in 2000 to 7.5\% in 2022. Quantitatively, the number of elderly people reached a staggering 373.7 million in 2022. This phenomenon, often referred to as the "Silver Tsunami," has had far-reaching socioeconomic and health impacts. In Malaysia, this trend is particularly pronounced; the population aged 60–64 years increased from 4.1\% to 6.0\%, while the population aged 65 and above surged to 6.1\% during the same period. This increase necessitates a repositioning of public health policies to provide infrastructure that adapts to the specific needs of older adults \cite{ASEANstats2023}.

Along with physical aging, the burden of mental health, particularly depression, is a crucial challenge threatening the well-being of older adults. A recent meta-analysis using a random-effects model estimated the prevalence of depression in the global elderly population at 19.2\% \cite{Jalali2024}. In the Malaysian context, this challenge is exacerbated by post-retirement social isolation. Research shows a strong linear correlation between equality and depression, with 63.9\% of retirees reporting high levels of equality and 64.8\% exhibiting symptoms of clinical depression \cite{Subramaniam2024}. Depression in older adults is not simply a mood disorder, but a systemic condition that impairs quality of life multidimensionally, encompassing physical, mental, social, and environmental aspects \cite{Park2022}. Physically, depression triggers a somatic cascade of chronic pain and persistent sleep disturbances, while mentally, it results in loss of anhedonia (loss of interest), feelings of helplessness, and severe decline in self-esteem \cite{Wroblewska2021}.

A major challenge in treating depression in the elderly lies in the complexity of early detection. There is a significant morphological discrepancy between the emotional expressions of young and older adults. Young individuals tend to display clear affect, such as drooping corners of the lips and droopy eyes when sad. In contrast, older individuals often exhibit a "flat expression" with a blank expression and a persistent furrowed brow. This is due to age-related declines in facial muscle elasticity and neuromuscular function, which are often used as conventional emotional indicators \cite{Grondhuis2021}. Although emotional expressions in the elderly are more subtle, advances in advanced facial analysis technology, particularly the observation of micro-expressions, offer promising new parameters for identifying signs of depression hidden behind physical facial changes \cite{KimH2023}.

In the clinical landscape, depression detection methodologies are polarized into two main approaches. Invasive approaches focus on analyzing biological biomarkers directly from the central nervous system. An example is the measurement of inflammatory cytokine levels such as interleukin-6 (IL-6) and interleukin-8 (IL-8) through lumbar puncture of cerebrospinal fluid (CSF) \cite{Sorensen2023}. Furthermore, Deep Brain Stimulation (DBS) techniques are used to disrupt neural circuits involved in the pathophysiology of depression through the placement of intracranial electrodes \cite{Figee2022}. While accurate, these methods carry high medical risks, are expensive, and pose significant discomfort for vulnerable elderly patients.

Alternatively, non-invasive approaches are becoming a primary choice due to their safety and minimal trauma. Psychometric instruments such as the Patient Health Questionnaire-9 \cite{Kroenke2001} and the Beck Depression Inventory \cite{Martin2010} have long been used to quantify symptom severity through patient self-report. However, these methods have limitations such as subjectivity and memory bias. Therefore, integrating clinical observation with non-verbal analysis based on the Facial Action Coding System (FACS) is highly relevant. Through FACS, microscopic changes in lip angles, brow tension, and eye blink frequency can be objectively coded as indicators of depression. The development of a non-invasive, automated, and accurate facial feature-based early detection system is urgently needed to ensure timely intervention and mitigate the long-term impact of depression on the physical and mental health of older adults in Malaysia and the ASEAN region at large.



\section{Problem Statement}
% Lorem ipsum dolor sit amet, consectetur adipiscing elit. Donec pulvinar pellentesque mauris in vulputate. Nam a porttitor tortor, ac lobortis sapien. Donec posuere nec libero et ornare. Morbi finibus mi nisl, in iaculis massa gravida ac. Ut quis nisi lorem. Donec luctus viverra velit at interdum. Nullam molestie sed ex eu commodo. Aenean volutpat eu risus at varius. Nunc finibus turpis id lectus tincidunt, nec porttitor felis semper. Phasellus consequat magna quis odio porta laoreet. Mauris non fringilla orci, eu auctor velit. Sed consequat sem nec pretium placerat. Vivamus in tortor venenatis erat tristique mattis vitae id elit. Donec pulvinar vulputate ante at fermentum. Integer scelerisque urna eget porttitor iaculis.
Depression is a complex mental health disorder and its prevalence is increasing, especially among the elderly population. This condition is not just an ordinary mood fluctuation, but a serious psychological problem that can systematically disrupt the functioning of daily life. Based on a study conducted by \cite{Devita2022}, the effects of depression in the elderly are very broad, including a drastic reduction in quality of life, disruption of sleep patterns like insomnia, to a decline in cognitive function such as memory degradation. Furthermore, depression has a close correlation with the exacerbation of chronic physical diseases, including cardiovascular disease, diabetes mellitus, and cancer. If there is no immediate intervention, depression in the elderly can reach a culmination point that can lead to death, namely the emergence of ideas and even the act of committing suicide. Given its massive implications, early detection is a medical urgency that needs to be researched more deeply through various scientific disciplines, especially the integration between health sciences and artificial intelligence technology.

Conventionally, the diagnosis of depression in the elderly relies heavily on subjective clinical methodology using self-administered questionnaires and face-to-face interviews. Commonly used instruments include the Geriatric Depression Scale (GDS-15), the Patient Health Questionnaire-9 (PHQ-9) \cite{Kroenke2001}, and the Beck Depression Inventory (BDI-II) \cite{Martin2010}. While these methods have served as the gold standard for many years, they have significant limitations. The effectiveness of these questionnaires depends heavily on honesty, emotional stability during completion, and memory capacity to recall symptoms. This poses significant challenges for older adults experiencing cognitive decline or those with verbal communication difficulties, often resulting in inaccurate or biased mental state outcomes.

In response to these limitations, the use of facial recognition technology has emerged as an innovative solution for more objective early detection of depression. The human face is a window into emotion, capable of depicting a person's psychological state through their microexpressions. However, applying this technology to older adults is challenging. As noted by Kumar Rao B et al. (2023), elderly faces have unique morphological characteristics, such as changes in skin texture like wrinkles, relaxed facial muscles, and dynamic changes in expression that differ significantly from those of younger populations. This unique structure demands an algorithm that is more adaptive and sensitive to the specific facial features of the elderly.

The technical basis of this research refers to a previous study by \cite{Sugiyanto2024}, which classified depression using 14 Action Unit (AU) features through a CNN-Pooling Less architecture. The study utilized the DAIC-WOZ and CASME-II datasets, validated with PHQ-8 scores, and used OpenSpace to render the face mesh. While achieving high accuracy on certain datasets (0.988 on DAIC-WOZ), the model's performance still fluctuated on other datasets (0.765 on CASME-II), indicating significant room for optimization, particularly in capturing the temporal dependencies of dynamic facial expressions.

Therefore, this study proposes the development of a more robust depression classification system by integrating a Convolutional Neural Networks-Bidirectional Long Short-Term Memory (CNN-BiLSTM) architecture. In this approach, the facial feature extraction technique will be enhanced using the MediaPipe instrument, which is capable of producing precise masking maps of up to 468 facial landmark points (Thaman et al., 2022). The selection of the CNN-BiLSTM hybrid model is based on functional synergy: CNN plays a role in extracting spatial features and visual patterns from images (Xie, 2023), while BiLSTM processes sequential data to learn bidirectional temporal dependencies. This combination allows the system to not only recognize static expressions but also more accurately capture the dynamics of emotional changes over time \cite{Febrian2022}.

Another fundamental challenge identified in this study is the scarcity of medically validated elderly facial datasets for depression labeling. The majority of currently available datasets are dominated by the productive age group. To address this data gap, this study applies a Transfer Learning strategy. The process is carried out in two stages: first, the model will be pre-trained using a clinically validated general facial dataset. Second, the model will undergo a fine-tuning process using an elderly facial dataset to adapt the algorithm's weights to the unique demographic features of older adults. Through this two-stage training scheme, this research is expected to produce a depression detection system with high reliability, superior accuracy, and the ability to serve as a valid diagnostic tool for medical professionals in managing the mental health of older adults.

% \section{Objective of Research and Scope of Works}
% \label{sec:examples}

% The first few paragraphs of Lorem Ipsum are given below.

\section{Research Question}
Research Question (labelled as RQ) that will be evaluated in this research as per below:

\begin{description}[leftmargin=1.2cm, labelsep=0.2cm]
    \item[RQ1 :] How does the extraction of facial landmarks from medically validated and synthetically aged datasets facilitate effective feature representation for two-phase Transfer Learning in classifying depressive emotional states in older adults?
    
    \item[RQ2 :] To what extent does the two-phase Transfer Learning strategy pre-training on expert-validated data and fine-tuning on older adult datasets—enhance the performance of CNN-BiLSTM models in detecting depression compared to single-phase training approaches?
\end{description}

\section{Research Hypothesis}
The main hypothesis in this study as per below :

\begin{enumerate}[label=(\arabic*), leftmargin=1.2cm]
    \item The CNN-BiLSTM model trained using a two-phase Transfer Learning approach—pre-training on medically validated depression datasets followed by fine-tuning on older adult facial datasets will achieve significantly higher classification accuracy and F1-score in detecting depression among older adults compared to models trained using only a single-phase approach.
    
    \item The adaptation of the CNN-BiLSTM model through fine-tuning on older adult facial data will result in improved model generalization and robustness, as measured by cross-dataset validation performance, due to the model enhanced ability to capture age-specific facial landmark features relevant to depression detection.
\end{enumerate}

\section{Research Scope}

This interdisciplinary research combines psychology and technology through the development of a deep learning-based system for early detection of depression in the elderly using facial analysis. The study's primary focus is on extracting facial marker features using the MediaPipe framework from two dataset categories: clinically validated facial images for the pre-training phase and elderly facial images for the fine-tuning phase. By implementing a CNN-BiLSTM architecture, this research aims to build a precise classification model that optimally adapts to the unique facial morphological characteristics of the elderly population.

A two-phase transfer learning strategy is the core innovation of this research, addressing the scarcity of medically validated facial image data for the elderly. The study's scope is limited to technical development and model performance evaluation using standard classification metrics, without expanding to real-time clinical implementation or the use of non-facial biomarkers. These findings are expected to significantly contribute to the development of non-invasive, automated depression screening instruments and serve as a methodological reference for similar studies facing limitations in annotated data in specific populations.



% \subsubsection{Some Notes}

% Duis sed rutrum eros, quis tempus risus. Etiam pellentesque nisi odio, eget dignissim eros ultrices et. Aliquam leo massa, fermentum vel odio sed, ullamcorper molestie lorem.

% \subsubsection{And Further}
% Duis sed rutrum eros, quis tempus risus. Etiam pellentesque nisi odio, eget dignissim eros ultrices et. Aliquam leo massa, fermentum vel odio sed, ullamcorper molestie lorem.


% \subsection{Least-Squares with Forgetting Factor AdaptiveLaw}

% \section{Mechatronic Suspension System; History and a Brief Background}
% Nulla ipsum augue, feugiat ac laoreet quis, pretium ut magna. Class aptent taciti sociosqu ad litora torquent per conubia nostra, per inceptos himenaeos. Integer blandit placerat dictum.



% \section{Summary}
% Nulla ipsum augue, feugiat ac laoreet quis, pretium ut magna. Class aptent taciti sociosqu ad litora torquent per conubia nostra, per inceptos himenaeos. Integer blandit placerat dictum.

% Sed dolor justo, scelerisque sed rutrum quis, porttitor a mauris. Cras non auctor felis, rutrum fringilla risus. Integer at convallis erat, sit amet luctus turpis. Duis sed rutrum eros, quis tempus risus. Etiam pellentesque nisi odio, eget dignissim eros ultrices et. Aliquam leo massa, fermentum vel odio sed, ullamcorper molestie lorem. Integer lorem felis, adipiscing sit amet interdum eget, auctor at lorem. Aliquam ultricies tortor eu nibh facilisis tincidunt.
\input{chap-fibonacci}
\input{chap-goldenratio}



% references are listed in refs.bib
\bibliography{refs}

\appendix
% Each appendix chapter from a separate file
\input{app-details}
\input{app-code}
\end{document}
